\section{Jismning og`irlik kuchi ta'sirida vertikal harakati}

\begin{enumerate}

\item Yerdan 12 m balandlikdan boshlang‘ich tezliksiz jism tushadi. Tushish boshidan 1 s o‘tgach jism qanday balandlikda bo‘ladi? $g=10 m/s^2$.
\begin{enumerate}[label=\Alph*)]
\item 5
\item 7
\item 10
\item 12
\end{enumerate}

\item Erkin tushayotgan jismning o‘ninchi sekundda o‘tgan yo‘li, undan oldingi sekundda o‘tgan yo‘lidan necha metrga ko‘p ekanligini aniqlang. Jismning boshlang‘ich tezligi nolga teng. $g=10 m/s^2$.
\begin{enumerate}[label=\Alph*)]
\item 20
\item 5
\item 10
\item 15
\end{enumerate}

\item Otish momentidan jismning yerga tushish momentigacha bo‘lgan vaqt 3 s ga teng bo‘lishi uchun, uni yer sirtidan qanday tezlik bilan vertikal yuqoriga otish kerak? $g=10 m/s^2$.
\begin{enumerate}[label=\Alph*)]
\item 15
\item 25
\item 10
\item 20
\end{enumerate}

\item Jism yer sirtidan 20 m/s tezlik bilan vertikal yuqoriga otildi. U qanday maksimal balandlikka ko‘tariladi? $g=10 m/s^2$.
\begin{enumerate}[label=\Alph*)]
\item 10
\item 15
\item 25
\item 20
\end{enumerate}

\item Vertikal yuqoriga otilgan sharcha otilish nuqtasiga 2,4 s dan so‘ng qaytib keldi. Sharcha qanday (sm da) balandlikka ko‘tarilgan? $g=10 m/s^2$.
\begin{enumerate}[label=\Alph*)]
\item 360
\item 720
\item 144
\item 72
\end{enumerate}

\item Yerdan ikki jism vertikal yuqoriga otildi. Birinchi jismning boshlang‘ich tezligi ikkinchisinining boshlang‘ich tezligidan 4 marta katta. Birinchi jism ikkinchisiga qaraganda necha marta balandroq ko‘tariladi?
\begin{enumerate}[label=\Alph*)]
\item 4
\item 8
\item 16
\item 32
\end{enumerate}

\item Qaysidir vaqt momentida erkin tushayotgan jismning tezligi 6 m/s bo‘ladi. 2 s dan keyin jismning tezligi qanday bo‘ladi? $g=10 m/s^2$.
\begin{enumerate}[label=\Alph*)]
\item 26
\item 16
\item 36
\item 6
\end{enumerate}

\item Koptok qandaydir balandlikdan 4 m/s tezlik bilan vertikal pastga otildi. Harakatning dastlabki ikki sekundida koptokning o‘rtacha tezligini toping. $g=10 m/s^2$.
\begin{enumerate}[label=\Alph*)]
\item 12
\item 14
\item 16
\item 10
\end{enumerate}

\item Tosh 2 m/s tezlik bilan vertikal pastga otildi. Otishdan 1 s o‘tgach, toshning tezligi necha marta ortadi? $g=10 m/s^2$.
\begin{enumerate}[label=\Alph*)]
\item 2
\item 10
\item 6
\item 5
\end{enumerate}

\item Vertikal pastga otilgan jismning tezligi bir sekunddan so‘ng 6 martaga ortdi. Otishdan ikki sekund o‘tgach uning tezligi boshlang‘ich bilan solishtirganda necha marta ortadi?
\begin{enumerate}[label=\Alph*)]
\item 11
\item 6
\item 22
\item 15
\end{enumerate}

\item Agar ko‘tarilishda jismning tezligi otishdan 0,8 s o‘tgach ikki marta kamaygan bo‘lsa, u qanday tezlik bilan vertikal yuqoriga otilgan? $g=10 m/s^2$.
\begin{enumerate}[label=\Alph*)]
\item 8
\item 20
\item 12
\item 16
\end{enumerate}

\item 20 m balandlikdan doskaga tushgan metal sharcha undan 25 \% tezligini yo‘qotib sakraydi. Sharcha urilishdan necha sekund o‘tib doskaga ikkinchi bor tushadi? $g=10 m/s^2$.
\begin{enumerate}[label=\Alph*)]
\item 2
\item 3
\item 5
\item 1
\end{enumerate}

\item Tosh 50 m/s tezlik bilan vertikal yuqoriga otildi. Necha sekund o‘tgach uning tezligi 30 m/s bo‘lib pastga yo‘nalgan bo‘ladi? $g=10 m/s^2$.
\begin{enumerate}[label=\Alph*)]
\item 10
\item 5
\item 8
\item 12
\end{enumerate}

\item Jism 40 m balandlikdan 5 m/s boshlang‘ich tezlik bilan vertikal yuqoriga otildi. 2 s o‘tgach jism qanday balandlikda bo‘ladi? $g=10 m/s^2$.
\begin{enumerate}[label=\Alph*)]
\item 30
\item 40
\item 20
\item 10
\end{enumerate}

\item Jism vertikal yuqoriga 20 m/s tezlik bilan otildi. Jismning harakat boshidan 3 s ichida o‘tgan yo‘lini toping. $g=10 m/s^2$.
\begin{enumerate}[label=\Alph*)]
\item 15
\item 20
\item 30
\item 25
\end{enumerate}

\item Yerdan 8 m balandda bo‘lgan nuqtadan vertikal yuqoriga otilgan jism otish paytidan 2 s o‘tgach yerga tushadi. Jism qanday tezlik bilan otilgan? $g=10 m/s^2$.
\begin{enumerate}[label=\Alph*)]
\item 12
\item 6
\item 8
\item 10
\end{enumerate}

\item Jism 15 m balandlikdan 10 m/s boshlang‘ich tezlik bilan vertikal yuqoriga otildi. Necha sekunddan so‘ng u yerga tushadi? $g=10 m/s^2$.
\begin{enumerate}[label=\Alph*)]
\item 5
\item 4
\item 3
\item 2
\end{enumerate}

\item To‘p 2,4 m balandlikdan 1 m/s tezlik bilan vertikal pastga otildi. Yerga tushish paytida uning tezligi qancha bo‘ladi? $g=10 m/s^2$.
\begin{enumerate}[label=\Alph*)]
\item 7
\item 1
\item 14
\item 10
\end{enumerate}

\item To‘p 2,4 m balandlikdan 1 m/s tezlik bilan vertikal pastga otildi. Qancha vaqtdan so‘ng (ms da) u yer sirtiga yetadi? $g=10 m/s^2$.
\begin{enumerate}[label=\Alph*)]
\item 240
\item 100
\item 400
\item 600
\end{enumerate}

\item Jism 80 m balandlikdan erkin tushadi. U tushish vaqtining so‘nggi sekundida qanday yo‘l o‘tadi? $g=10 m/s^2$.
\begin{enumerate}[label=\Alph*)]
\item 15
\item 35
\item 25
\item 45
\end{enumerate}

\item Agar jism tushish vaqtining oxirgi sekundida 45 m yo‘l o‘tgan bo‘lsa, u qanday balandlikdan tushgan? $g=10 m/s^2$.
\begin{enumerate}[label=\Alph*)]
\item 100
\item 150
\item 125
\item 175
\end{enumerate}

\item Jism yer sirtidan vertikal yuqoriga otildi. Jismning maksimal ko‘tarilish balandligining 8/9 qismini tashkil qiluvchi balandlikdagi tezligi boshlang‘ich tezlikdan necha marta kichik?
\begin{enumerate}[label=\Alph*)]
\item 3
\item 9
\item 6
\item 2
\end{enumerate}

\item Jism vertikal yuqoriga otiladi. Kuzatuvchi qayd etdiki, jism 75 m balandlikdagi nuqtada 2 s vaqt intervali bilan ikki marta bo‘ldi. Jismning boshlang‘ich tezligini toping. $g=10 m/s^2$.
\begin{enumerate}[label=\Alph*)]
\item 30
\item 50
\item 40
\item 60
\end{enumerate}

\item Vertikal yuqoriga otilgan tosh bitta balandlikda ikki marta bo‘ldi – harakat boshidan 0,8 va 1,5 s vaqt o‘tganda. Bu balandlik qanchaga teng? $g=10 m/s^2$.
\begin{enumerate}[label=\Alph*)]
\item 1
\item 2
\item 4
\item 6
\end{enumerate}

\item Agar jismning harakatning oxirgi sekundida o‘tgan yo‘li dastlabki sekundda o‘tgan yo‘lidan besh marta katta bo‘lsa, jism qanday balandlikdan boshlang‘ich tezliksiz tushgan? $g=10 m/s^2$.
\begin{enumerate}[label=\Alph*)]
\item 45
\item 15
\item 30
\item 90
\end{enumerate}

\item Erkin tushayotgan jism o‘z harakatining oxirgi 10 sekundida butun yo‘lning 3/4 ini o‘tadi. Jism boshlang‘ich tezliksiz tushgan balandlikni aniqlang. $g=10 m/s^2$.
\begin{enumerate}[label=\Alph*)]
\item 10
\item 50
\item 200
\item 100
\end{enumerate}



\end{enumerate}


\begin{table}[h!]
\centering
\caption{To‘g‘ri javoblar jadvali (1-28)}
\begin{tabular}{|c|c|c|c|c|c|c|c|c|c|}
\hline
№ & J. & № & J. & № & J. & № & J. & № & J. \\ \hline
1 & B & 7 & A & 13 & C & 19 & D & 25 & A \\ \hline
2 & C & 8 & B & 14 & A & 20 & B & 26 & C \\ \hline
3 & A & 9 & C & 15 & D & 21 & C &  &  \\ \hline
4 & D & 10 & A & 16 & B & 22 & A &  &  \\ \hline
5 & B & 11 & D & 17 & C & 23 & C & & \\ \hline
6 & C & 12 & B & 18 & A & 24 & D & & \\ \hline
\end{tabular}
\end{table}